%! The Statistical Cycle and Types of Data — Unit 1, Lesson 1.1 — Lesson Plan
% PILOT SHAPE (2026-09-06), Main Ideas / Notes table (2026-09-12): regenerated from the
% notes-only shape (5/30/10/7/3, seven notes sections, a We Do, a solo practicebox and an
% extensionbox — all retired). GRADUAL-RELEASE plan (warm-up / guided notes and practice /
% spoken debrief / close and START THE HOMEWORK). The guided notes carry the release ROW BY
% ROW in one Main Ideas / Notes table: I do / we do in rows 1-4 (problems 1-15) -> we do the
% Guided Practice row (problems 16-19). THERE IS NO INDEPENDENT PRACTICE SET IN THE NOTES —
% the homework is the individual practice, and the period ends with students starting it.
% Teacher-facing. Breakable boxes are `skillbox` with a \boxguard ahead; the two tabularx
% boxes are `fixedskillbox` (never splits). No group activity, no tier boxes, no exit ticket.
\documentclass[10pt]{article}
\usepackage{saar-article}
\usepackage{saar-boxes}
\usepackage{graphicx}   % -article does NOT load graphicx
\graphicspath{{images/}}
\newcommand{\TallMath}[1]{$\displaystyle #1\rule[-1.4em]{0pt}{3.2em}$}
\newcommand{\UnitNumberName}{Unit 1: Asking Questions with Data: Study Design \quad}
\newcommand{\LessonNumberName}{Lesson 1.1: The Statistical Cycle and Types of Data}

\begin{document}
\begin{center}
    {\Large\bfseries \CourseName \\
    {\small\bfseries \UnitNumberName \LessonNumberName}}
\end{center}

\begin{tcolorbox}[colback=frost, colframe=cerulean, boxrule=0.9pt, arc=2mm,
    left=2mm, right=2mm, top=1.5mm, bottom=1.5mm]
    \textbf{Primary Objective:} Students will name the four stages of the data cycle and say
    what each hands to the next, test a question against three checks and repair one that
    fails, decide from the \emph{question alone} whether the data will be categorical or
    quantitative, organize counts into a frequency table with relative frequencies, write a
    conclusion that names the group it applies to, and name the stage at which a flawed study
    broke down.
    \\[2pt]
    \textbf{Standards:} PS.DC.1a, PS.DC.1b, PS.DC.1c; AFDA.DA.2a \quad$\cdot$\quad previews
    PS.DC.1d--e (Lesson 1.2)
    \\[2pt]
    \textbf{Lesson model:} \emph{Gradual release --- I do, we do, you do --- all of it inside
    the guided notes.} There is no group activity. The notes name the vocabulary as they teach
    it and deliver the instruction \textbf{row by row in one Main Ideas / Notes table} (four
    rows, problems 1--15: you read the definition, mark up the display, and work the first
    problem of each grid; the class works the rest); the Guided Practice row (problems 16--19)
    is worked together with students holding the pen; \textbf{the homework is the individual
    practice}, and students start it in class while you circulate. The whole-class debrief is
    \emph{spoken} --- there is no practice set at the end of the notes and no exit ticket.
    \\[2pt]
    \textbf{Note on scope:} \emph{sample} and \emph{population} are used informally today
    (``the $60$ she asked'' versus ``all $900$''). Do not name them --- Lesson 1.2 does, and
    the homework's closing box is the hand-off.
\end{tcolorbox}

\renewcommand{\arraystretch}{1.4}
\boxguard
\begin{skillbox}[Learning Targets \& Key Understandings]{goldbox}
\begin{tabularx}{\linewidth}{@{}X|X@{}}
\textbf{Learning targets (student language)} & \textbf{Key understandings (the ``why'')} \\[2pt]
\hline
\vspace{-6pt}
\begin{itemize}[leftmargin=1.1em, nosep, after=\vspace{2pt}]
  \item I can name the four stages of the \textbf{data cycle} and say what each stage hands
        to the next. \hfill \textit{PS.DC.1a; AFDA.DA.2a}
  \item I can test a question against three checks --- answers \textbf{vary}, the
        \textbf{group} is named, the \textbf{measurement} is named --- repair one that fails,
        and name the \textbf{variable of interest}. \hfill \textit{PS.DC.1b; AFDA.DA.2a}
  \item I can decide from the \emph{question alone} whether the data will be
        \textbf{categorical} or \textbf{quantitative}, and build a \textbf{frequency table}
        with \textbf{relative frequencies}. \hfill \textit{PS.DC.1c; PS.DC.1a}
  \item I can write a \textbf{conclusion in context}, say what the data do \emph{not} allow,
        and name the stage a flawed study broke at. \hfill \textit{PS.DC.1a--b}
\end{itemize}
&
\vspace{-6pt}
\begin{itemize}[leftmargin=1.1em, nosep, after=\vspace{2pt}]
  \item \emph{The stages are ordered because each one consumes the previous one's output.} A
        conclusion cannot name a group stage 1 never named.
  \item \textbf{Target misconception:} that the \emph{data} decide the type --- so you must
        collect before you can classify. The \textbf{question} decides it, before anything
        exists: \emph{``for each one I will record \underline{\hspace{1.2cm}}''} settles it.
        Lesson 1.0 classified variables sitting in a finished table; today the decision comes
        first.
  \item \textbf{Second misconception:} that a study broke wherever the bad \emph{sentence}
        appeared --- always stage 4. Ask where the problem \emph{entered}. Problem~14's
        arithmetic is correct and the study broke at \textbf{stage 1}.
  \item Stage 3 adds no information; it makes information \emph{visible}. Sixty answers in a
        list and sixty in a frequency table are the same data --- only one can be read.
  \item A conclusion may answer \textbf{only} the question that was asked, about
        \textbf{only} the group that was asked. Change either and the same correct percent
        supports nothing.
\end{itemize}
\end{tabularx}
\end{skillbox}

\boxguard[24]
\begin{skillbox}[{Vocabulary, Concepts, \& Theorems --- taught directly in the guided notes}]{greenbox}
\small
\renewcommand{\arraystretch}{1.35}
\begin{tabularx}{\linewidth}{@{} >{\bfseries}p{3.1cm} X @{}}
Data cycle & The four repeating stages of a study: \emph{formulate questions $\rightarrow$
       collect or acquire data $\rightarrow$ organize and represent $\rightarrow$ analyze and
       communicate}. Also called the statistical cycle. \\
Statistical question & A question whose answers are expected to vary, that names the group it
       is about, and that says what will be recorded for each individual. \\
Variable of interest & The characteristic the question says you will record for each
       individual. \\
Categorical / quantitative & A variable whose values are labels / a variable whose values are
       amounts arithmetic means something on. \textbf{Decided by the question.} \\
Frequency & The count of individuals in one category or at one value. \\
Frequency table & A table pairing each category or value with its frequency. \\
Relative frequency & A count divided by the total, written as a decimal or a percent --- the
       share of the whole that group represents. \\
Conclusion in context & A sentence stating what the data show, naming the group it applies to
       and what was measured. \\
\end{tabularx}
\end{skillbox}

% fixedskillbox never splits, so the phase table stays intact on one page.
% THE PHASE TABLE IS A CONTRACT: 5 / 35 / 10 / 10 — the I do is ~20 of the 35 and the Guided
% Practice ~15; the debrief is spoken; the homework start is the second formative read.
\begin{fixedskillbox}[Lesson at a Glance --- \MeetingLength]{frost}
\renewcommand{\arraystretch}{1.25}
\begin{tabularx}{\linewidth}{@{}l c X X@{}}
\textbf{Phase} & \textbf{Min} & \textbf{Students} & \textbf{Teacher} \\
\hline
Warm-Up & 5 & The council's bus percent, the scale-up, then two variables to classify ---
  silent & Debrief item~3 aloud; write ``$60$ asked, $900$ reported on'' on the board and
  leave it up all period \\
Guided Notes \& Practice & 35 & Fill the vocabulary box as each term is named; work rows
  1--4 of the table (problems 1--15) with the pen in their hand; then the Guided Practice row,
  \emph{The Cycle Turns} (16--19), together & \textbf{I do / we do, row by row}: read the
  definition, mark up the display, work problems 1, 5, 9's bus row, 12; the class works the
  rest (20) $\to$ \textbf{we do} Guided Practice 16--19 (15) \\
Debrief & 10 & Pens down, papers up --- answer aloud, nothing to fill in & Put the council's
  ``$10\%$ bike now'' claim back on the board and take the ``the arithmetic is right so the
  claim is right'' reflex, then the ``stage 4'' reflex \\
Close \& Assign & 10 & \textbf{Start the homework in class}, alone and silent & Announce the
  homework and its form (packet or Desmos), circulate for the second formative read, preview
  Lesson~1.2 \\
\end{tabularx}
\end{fixedskillbox}

\begin{fixedskillbox}[Warm-Up --- Activate Prior Knowledge (5 min)]{frost}
\begin{tabularx}{\linewidth}{@{}X X@{}}
\begin{minipage}[t]{\linewidth}\vspace{0pt}
    \textbf{The three items and what each seeds}
    \begin{itemize}[leftmargin=1.1em, itemsep=1pt, topsep=2pt]
      \item \textbf{1.} $27$ of $60$ as a percent. \textbf{Seeds:} the bus row of the
            frequency table in row~3 --- problem~9 hands it to them already done, so the
            table starts from their own arithmetic.
      \item \textbf{2.} $45\%$ of $900$. \textbf{Seeds:} the scale-up problem~11 asks for and
            Guided Practice problem~17 asks again.
      \item \textbf{3.} Two variables, classified from a finished study (Lesson~1.0).
            \textbf{Seeds --- and is taken away by --- row~4:} today the \emph{question}
            decides the type, before any data exist.
    \end{itemize}
\end{minipage}
&
\begin{minipage}[t]{\linewidth}\vspace{0pt}
    \textbf{Running it}
    \begin{itemize}[leftmargin=1.1em, itemsep=1pt, topsep=2pt]
      \item \textbf{Debrief item~3 aloud} and write the two counts on the board --- \emph{$60$
            asked, $900$ reported on} --- in a student's own words. Leave it up all period;
            row~3 scales from that line and the hook is settled against it.
      \item Item~3's reason column is the tell. Accept \emph{``bus, car, walk, bike are
            labels''}; reject \emph{``because it's words.''} The reason must be that
            arithmetic on those values means nothing.
      \item Items 1--2 are arithmetic they already own --- say so, so nobody thinks they are
            being retested. Do not let them expand.
    \end{itemize}
\end{minipage}
\end{tabularx}
\end{fixedskillbox}

\boxguard
\begin{skillbox}[Hook --- before anyone sits down]{frost}
The hook is on \textbf{page 1 of the notes}, under the vocabulary box, and on the board:
\textit{``The Student Council wants more bike racks. Last spring they asked $60$ students how
they usually get to school; $6$ said bike. Their claim: $10\%$ bike now, so a lot more would
if there were somewhere to park.''} \textbf{Say the arithmetic is right} --- $6$ of $60$ really
is $10\%$ --- then ask: \textbf{does the survey support the claim?} Students circle
\emph{yes} / \emph{no} / \emph{can't tell} and write one line of reason in the hook box
\emph{before} anyone speaks; then take the hand vote and write the split on the board. Most
rooms split, and a large group votes \emph{yes} because the percent checks out.
\textbf{Do not settle it.} Say only that the answer is problem~14, in \emph{The Question
Decides} row, and that they will vote again there. The vote is what makes the crux land:
students have to have committed to \emph{correct arithmetic means a correct claim} before the
notes take it away.
\end{skillbox}

\boxguard
\begin{skillbox}[Guided Notes \& Practice --- I do, then we do (35 min)]{frost}
\textbf{This is the whole lesson.} There is no group activity, and there is no independent
practice set on this page: \textbf{the homework is the individual practice}, started in class.
The notes are \textbf{one Main Ideas / Notes table}: four instruction rows (problems 1--15),
then the Guided Practice row (16--19). \textbf{The rhythm in every row is the same}: read the
definition aloud, mark up the display, work the first problem of the grid yourself, then hand
the rest to the room with the pen in their hand. The vocabulary box on page~1 is filled
\emph{as each term is named}, never in advance. The only blanks are the frequency table of
problem~9; every other answer is written in a problem's own space.
\begin{multicols}{2}
\textbf{I do / we do, row by row (about 20 min).}

\emph{Row 1, Data Cycle (problems 1--4).} Name the four stages, then \textbf{run the third
column of the picture} --- what each stage \emph{hands to the next}. Ask: \emph{if stage 1
never says who the study is about, what can stage 4 say?} Nothing, about anybody. Then the
grid is the council's own four actions out of order: work problem~1 aloud (stage~2), and take
2--4 from the class ($4$, $1$, $3$). Four minutes.

\emph{Row 2, Three Tests (problems 5--8).} Vary / group / record. Lesson~1.0 gave test~1
only, so 2 and 3 are the new material. Work problem~5 yourself (``How do I get to school?''
fails \emph{vary}). The class takes 6 (fails \emph{group}) and 7 (fails \emph{record} ---
``good'' is not something you write down). \textbf{Problem~8 is the trap: it passes all three
as written.} Take a hand count before anyone writes --- most rooms start rewriting it because
the task said repair. \textbf{Insist a repair keeps the original subject.} Five minutes.

\emph{Row 3, Frequency Table (problems 9--11).} The bus row is done for them from the warm-up
--- do that row aloud, then hand the other three to the room. Check the percents sum to
$100\%$. \textbf{Ask: what did stage 3 add?} Nothing --- it made the counts readable. Say it
out loud; students treat organizing as busywork. Problem~11 is the scale-up and the first
conclusion sentence. Four minutes.

\emph{Row 4, The Question Decides (problems 12--15) --- the crux.} Mark up the display: the
\emph{same} $30$ students, two questions, two types. \textbf{Ask: do you have to collect the
data to find out whether they are categorical?} No. Work problem~12 yourself; the class takes
13, then the \emph{In your own words} line. Then \textbf{re-take the hook vote at problem~14
before you say anything}: the arithmetic is right, and the survey still cannot support
\emph{would bike} --- it asked how students travel \emph{now}, so the problem entered at
\textbf{stage 1}. Problem~15 is the counterweight: the cart's $35\%$ is also correct, and that
one entered at \textbf{stage 2}. Seven minutes.

\columnbreak
\textbf{We do (about 15 min).} The Guided Practice row, \emph{The Cycle Turns} (problems
16--19) --- the council goes back to stage 1 with the question the first survey could not
answer, and the row runs the same moves the homework will ask alone.
\textbf{Students hold the pen; you hold the questions.} Problem~16 goes on them cold --- the
variable and its type, \emph{before} anyone looks at the $80$ answers. Problem~17 is warm-up
arithmetic ($30\%$, about $270$ of $900$). Problem~18 is the conclusion sentence: make them
supply the group, the number, and the measurement separately. \textbf{Problem~19 is the crux
transferred}: if the $80$ had all been standing at the bike rack, the problem enters at
\textbf{stage 2}, and the report can no longer describe Riverbend students at all.

Circulate during Guided Practice and demand the reason, not the word:
\begin{itemize}[leftmargin=1.1em, nosep]
  \item \textbf{Problem 16:} ``What will you write down for each student? Then what type is
        it?'' (Not: \emph{what do the answers look like} --- there are none yet.)
  \item \textbf{Problem 17:} ``Who does the $270$ describe --- the $80$, or the $900$?''
  \item \textbf{Problem 18:} ``Which $80$ students? Asked what, exactly?''
  \item \textbf{Problem 19:} ``\textbf{Where did the problem \emph{enter}?} Who was never
        asked?'' A student who writes \emph{stage 4} has named where the sentence appeared.
\end{itemize}
While circulating, pick the papers to put up in the debrief --- one problem~19 that said
\emph{stage 2} with a reason, and one that said \emph{stage 4}. \textbf{Your formative read
comes twice}: here, and again in the last ten minutes as students start the homework.
\end{multicols}
\end{skillbox}

\boxguard[24]
\begin{skillbox}[Debrief --- whole class, spoken (10 min)]{frost}
Pens down, papers up. \textbf{This debrief is spoken} --- there is no box at the end of the
notes to fill in, so it lives entirely on the board and in the room.
\begin{multicols}{2}
\begin{enumerate}[leftmargin=1.3em, nosep, itemsep=3pt]
  \item \textbf{Put the council's claim back up} and make a student say the sentence: the
        arithmetic is right, and the survey still does not support it, because stage~1 asked
        how students travel \emph{now}. Then state it as PS.DC.1b: a conclusion answers only
        the question that was asked.
  \item Put up the two problem~19s you picked. The difference between them is not
        arithmetic --- it is whether the student names where the problem \emph{entered} or
        where the sentence appeared. Read both aloud.
  \item Put problem~12 back up and take the \emph{you have to see the data first} reflex:
        \emph{``Nothing has been collected. What will you write down for each student?''}
  \item Take problem~15 last --- the cart's correct $35\%$ --- and get \emph{both} unsupported
        claims: ``Riverbend students'' when only cart customers were asked, and ``every
        morning'' when one morning was watched.
\end{enumerate}
\columnbreak
\textbf{Check for understanding --- ask it cold, whole class.} \emph{``The yearbook staff
surveyed only the students in the yearbook class, then wrote: Riverbend seniors prefer song
lyrics in the yearbook. Which stage did the problem enter at, and why does that stop the
conclusion?''} Hands down, thirty seconds of think time, then cold-call three students.
Correct answer: \textbf{stage 2} --- the data describe only the class surveyed, so the
conclusion may not be extended to seniors generally. \emph{This replaces the exit ticket.}

\textbf{Formative read.} Sort the three answers as they come: \textbf{(a)} named stage~2 and
said why; \textbf{(b)} named stage~2 but could not say what it stops; \textbf{(c)} named
stage~4 because that is where the bad sentence is. Pile (b) is the teachable one. \textbf{If
pile (c) runs past a third of the room, open Lesson~1.2 by re-running problems 14--15}, and
say so plainly.

\textbf{If the debrief runs short}, cut items 3 and~4 --- item~1 is the one that cannot be
cut. \textbf{Do not let the debrief eat the last ten minutes} --- the homework start is not
optional padding, it is your second formative read.
\end{multicols}
\end{skillbox}

\boxguard
\begin{skillbox}[Homework --- scored, started in class, due after two study halls]{goldbox}
Six exercises on the \textbf{Lakeside Farmers Market} (about $800$ Saturday shoppers; $50$
asked one Saturday): \textbf{1} name the variable of interest and its type with a reason;
\textbf{2} the frequency table with relative frequencies (\emph{the spiral to Lesson 1.0});
\textbf{3} scale $10\%$ to $800$ and write the conclusion sentence; \textbf{4} the
\textbf{contrast pair} --- the same $50$ shoppers, two questions, two types, decided before
anything is collected; \textbf{5} assign a stage to four failures, one per stage;
\textbf{6} a report whose arithmetic is entirely correct and whose conclusion still does not
follow --- \emph{the crux head-on}. Capped at \textbf{two pages} (user decision 2026-09-06).

\textbf{Start it in the last ten minutes of the period, in class and alone.} \textbf{Scoring:}
12 points --- 2 per item. \textbf{Item~6 is where the misconception shows}: a student who
answers ``stage 4, that's where she wrote it'' has learned the word and not the rule.
\textbf{Item~4's last line carries the second check} --- what made the two types different is
\emph{the question}, not the shoppers.

\textbf{Packet or Desmos --- decided here.} The packet pages are authored and are the default.
Desmos Classroom has percent-of-a-total practice that would cover items 2--3, but nothing that
asks a student to classify from a question with no data, or to explain why a correct percent
does not support a sentence --- the items this lesson exists for --- so \textbf{1.1 is a
packet night}. Say so aloud at Close \& Assign.

\textbf{Preview of Lesson~1.2.} The manager asked $50$ shoppers and then wrote about all
$800$. Do not name those groups here; the homework's closing box is the hand-off and 1.2 opens
on it.

\textbf{Connections \& Big Ideas.} \textit{Unit 1 core idea:} a conclusion is only as
trustworthy as the study that produced it. \textit{Standards covered:} PS.DC.1a, PS.DC.1b,
PS.DC.1c; AFDA.DA.2a. \textit{Spirals forward to:} population vs.\ sample and parameter vs.\
statistic (1.2, PS.DC.1d--e), sampling technique (1.3, PS.DC.2a--b), bias (1.4, PS.DC.2c), and
the full-cycle project (1.8). \textit{Reaches back to:} Lesson 1.0's individuals, variables,
and variable types, and Algebra~1 percent and proportional reasoning.
\end{skillbox}

\boxguard
\begin{skillbox}[Watch For (while circulating)]{redbox}
\begin{itemize}[leftmargin=1.2em, nosep]
  \item \textbf{``You can't tell the type until you see the data''} (problems~12--13 and~16,
        HW~4 --- the target misconception). Probe: \emph{``Nothing has been collected. What
        will you write down for each one?''}
  \item \textbf{``Stage 4, because that's where they said it''} (problems~14, 15 and~19,
        HW~5--6). Probe: \emph{``Where did the problem \emph{enter}? What was never asked?''}
  \item \textbf{``The arithmetic is right, so the claim is right''} (problem~14, the hook
        vote, HW~6). Probe: \emph{``Read the question they asked. Does it contain the word
        \emph{would}?''}
  \item \textbf{Repairs that change the topic} (problems~5--7, HW~4). Students rewrite a
        brand-new question. Probe: \emph{``Keep her subject. Put back only what was
        missing.''} And problem~8 needs no repair at all.
  \item \textbf{Percent-of-what} (problem~9, HW~2): dividing by the largest count instead of
        the total. Probe: \emph{``Add your four percents. What should they come to?''}
  \item \textbf{A conclusion with no group in it} (problems~11 and~18, HW~3): ``$30\%$ would
        ride.'' Probe: \emph{``Thirty percent of whom?''}
\end{itemize}
Cold-call prompts: \emph{``What does stage 2 hand to stage 3?''} \quad \emph{``Name the
variable of interest in that question.''} \quad \emph{``Say what this conclusion is not
allowed to claim.''}
\end{skillbox}

\boxguard
\begin{skillbox}[Close \& Assign (10 min)]{goldbox}
\textbf{Students start the homework here, in class, alone and silent --- this is the individual
practice, and the first minutes of it are your best formative read.} Hand the launch first ---
six exercises on paper, scored out of 12, due at the start of the first class after two study
halls --- and point out that the \emph{Keep in Mind} box on the cover is the page to flip back
to. Circulate: \textbf{item~6 is the column that tells you whether \emph{The Question Decides}
row landed}, and it is on page~2, so put students on 1, 4, and~6 first and let 2, 3, and~5 go
home. Sort what you see into three piles as you walk --- named the stage the problem entered
at and said what it stops (ready); named the right stage with no reason (ready, needs the
language); wrote \emph{stage 4} (the misconception survived) --- and open Lesson~1.2 by
re-running problems 14--15 if that third pile ran past a third of the room. Then close on what
changed: \emph{``You walked in able to turn $27$ of $60$ into a percent. You walk out knowing
that the question decides what kind of data you get before any exist, that a study runs in
four stages and each one only gets what the last one handed over --- and that a percent can be
perfectly correct and still not let you say the sentence you wanted to say.''}
\textbf{Preview:} Lesson~1.2, \emph{Populations, Samples, Parameters, and Statistics} --- the
council asked $60$ students and then talked about all $900$. Those two groups have names, and
so do the numbers computed from each.
\end{skillbox}

% ── Teacher notes — one per component, in packet order (three) ──────────────
% Teacher-only prose lives HERE, never in a _key: a note is the one block in a key with no
% counterpart in the blank. No note for the debrief or the homework start — both are fully
% specified in their own boxes above.
\begin{teachernote}[Warm-Up]
Five minutes, silent, timed. Items 1 and~2 are arithmetic students already own --- say so, so
nobody thinks they are being retested. \textbf{Item~3 is the one to read while collecting}:
expect the classification to be right and the \emph{reason} column to be thin. Accept
``bus, car, walk, bike are labels'' but not ``because it's words''; the reason has to be that
no arithmetic on those values means anything. Note who classified correctly \emph{here} ---
row~4 takes the finished table away and makes them decide from the question instead, and the
students who leaned on the visible values are the ones who will stall at problem~12. Write the
two counts on the board in a student's own words and \textbf{leave them up all period}; row~3
scales from that line and the debrief circles it again. Do \emph{not} supply the words
\emph{sample} and \emph{population} --- they belong to Lesson~1.2.
\end{teachernote}

\begin{teachernote}[Guided Notes \& Practice]
\textbf{This one document is the lesson --- pace it 20 / 15, then 10 for the debrief, then 10
for the homework start.} It is one table, and every row runs the same way: you read the
definition, mark up the display, and work the first problem of the grid (1, 5, 9's bus row,
12); the class works the rest with the pen in their hand. Rows 1 and~3 must move: four minutes
each, because the arithmetic is warm-up-familiar and the content is the third column of the
picture, not the boxes.

\textbf{Take the hand count on problem~8 before anyone writes.} It passes all three tests as
written, and the grid's own instruction says \emph{repair it} --- so the room starts
rewriting. A student who repairs a working question has learned to rewrite, not to test. Let
them commit, then run the three tests on it one at a time out loud.

\textbf{Spend the time in row~4, \emph{The Question Decides}}, which carries the
misconception twice over. Work problem~12 live --- nothing has been collected, and the type is
already decided --- then let the class do~13 and the \emph{In your own words} line, and
\textbf{re-take the hook vote at problem~14 before you say anything}: yes, no, or can't tell.
Let the room argue. Do not resolve it by repeating the definition; read the council's actual
question off the board and ask whether the word \emph{would} is anywhere in it. The arithmetic
being right is exactly what makes this land. Problem~15 is the counterweight: another correct
percent, a different entry stage (\textbf{2}, not 1), so ``where did it enter'' stays a real
question rather than a single memorized answer.

\textbf{Guided Practice (16--19) is where the pen changes hands.} \emph{The Cycle Turns} runs
the four moves the homework will ask alone --- variable and type before the data, the percent
and the scale-up, the conclusion sentence, the diagnosis --- on the council's second survey,
so the cycle visibly turns rather than being described. If you only get through 16, 17, and
19, that is the right trade; problem~18's sentence can be said aloud.

\textbf{There is no independent practice set on this page, and that is deliberate --- the
homework is the individual practice.} Guided Practice is the last row; the period ends with
students starting the homework in front of you, which is where your second formative read
comes from. Item~6 is the one to watch. If the \emph{stage 4} pile is more than a third of the
room, \textbf{open Lesson~1.2 by re-running problems 14--15}, and say so plainly.

\textbf{Debrief (10 min), spoken.} The council's claim back on the board with a student saying
why a correct percent does not support it; it is the one move that cannot be cut. \textbf{Do
not let the debrief eat the last ten minutes} --- the homework start is not optional padding.
\end{teachernote}

\begin{teachernote}[Homework]
\textbf{Scored, 12 points (2 per item), due at the start of the first class after two study
halls.} The ten minutes at the end of the period are a \emph{start}, not a completion: put
students on 1, 4, and~6 while you can still catch a wrong start, and let 2, 3, and~5 go home.
Item~6 is the one to read first when scoring --- any student who explains the bad conclusion
by ``she wrote it wrong in the report'' has named stage~4 and missed that nobody was ever
asked about a shuttle, which is a stage-1 failure. Item~4 is the fullest diagnostic in the
set: the same $50$ shoppers appear in both rows, so a student who says the type depends on the
answers has nowhere to hide. Score items 1--5 for correctness and item~6 for the \emph{reason},
not the verdict. \textbf{When you score the page, sort item~6 into three piles} --- named the
entry stage and what it stops (ready); right stage, no reason (ready, needs the language);
\emph{stage 4} (the misconception survived). If more than a third land in the last pile,
reopen the council's claim in the Lesson~1.2 warm-up debrief rather than letting it ride.

\textbf{Packet is the call for 1.1.} Desmos carries the percent arithmetic and nothing else in
this set; do not swap it in. Item~4's ``same shoppers, two questions'' and item~6's
correct-arithmetic / wrong-conclusion reasoning have no Desmos equivalent. Reopen the $50$
shoppers versus all $800$ at the start of Lesson~1.2 --- it is 1.2's opening.
\end{teachernote}
\end{document}
